This section describes how to use Claude Code as an AI assistant within the
task-based research workflow.
The approach draws on Chris Blattman's agentic workflow guidelines
(\url{https://claudeblattman.com/essentials/})
and adapts them for the Make-based pipeline architecture.

\subsection{The CLAUDE.md knowledge base}

Every project should have a \texttt{CLAUDE.md} file at the repository root.
This file is automatically loaded at the start of every Claude Code session
and serves as the AI's knowledge base about the project.

Key sections to include:
\begin{itemize}
	\item Project overview: what the project does, who works on it.
	\item Repository structure: where things live.
	\item Build and run instructions: how to execute tasks locally and on HPC.
	\item Task directory convention: the input/code/output structure.
	\item Code quality rules: what the CI/CD checks enforce.
	\item Commit conventions: message prefixes and formatting.
\end{itemize}

Keep \texttt{CLAUDE.md} under 200 lines.
Use \texttt{<!-- CUSTOMIZE -->} markers for sections that change per-project.
Test it by asking Claude ``What do you know about this project?''

\subsection{The Prompt-Plan-Review-Revise loop}

This is the core agentic workflow, adapted from Blattman's four-step process:

\begin{enumerate}
	\item \textbf{Prompt}: structure your thinking.
	Use \texttt{/prompt} to convert stream-of-consciousness into a structured prompt
	with Role, Context, Task, Constraints, Output Format, and Bookend sections.
	\item \textbf{Plan}: use plan mode (\texttt{Shift+Tab}).
	Claude reads files and designs an implementation approach,
	but cannot execute changes until you approve.
	For non-trivial tasks, assign multiple perspectives to surface hidden assumptions.
	\item \textbf{Review}: use \texttt{/review-plan} to spawn an independent agent.
	This fresh agent critiques the plan without conversational self-bias.
	It uses Red/Yellow/Green feedback to flag blocking issues, risks, and strengths.
	Run review multiple times---each round catches different problem categories.
	\item \textbf{Revise}: incorporate feedback and iterate.
	Only approve the plan when Red items are resolved.
\end{enumerate}

The key principle: \textbf{separate planning from reviewing.}
Claude reviewing its own plan in the same conversation is like asking someone
to peer-review their own paper.

\subsection{Custom skills}

The project includes custom Claude Code skills in \texttt{.claude/commands/}:

\begin{description}
	\item[\texttt{/review-plan}]
	Spawns a fresh agent to critique the current plan.
	Uses Red (blocking), Yellow (risks), Green (strengths) format.
	Provides a Verdict: Approve / Approve with Changes / Revise.
	\item[\texttt{/review-code}]
	Research-specific code review checking reproducibility,
	portability (local + DCC + TACC), pipeline integrity,
	and code quality.
	\item[\texttt{/done}]
	End-of-session wrap-up.
	Captures what was accomplished, decisions made, open questions,
	and writes a handoff note to \texttt{.claude/handoff.md} for next-session continuity.
	Generates a project manager update.
	\item[\texttt{/prompt}]
	Restructures stream-of-consciousness input into the six-section prompt format.
\end{description}

\subsection{Session continuity}

Research projects span many sessions.
Use \texttt{/done} at the end of each session to write a handoff note.
The next session can read \texttt{.claude/handoff.md} to understand where things left off.

The project manager update template (from \texttt{CLAUDE.md}):
\begin{lstlisting}[language=bash]
**PROJECT update**: <1-3 sentence summary>

**Next steps**:
1. <step>
2. <step>
\end{lstlisting}

\subsection{Design principles for agentic research}

From Blattman's guidelines, adapted for our workflow:

\begin{enumerate}
	\item \textbf{One skill, one job.}
	Focused skills outperform multi-purpose ones.
	\item \textbf{Default to useful action.}
	Skills should work with no arguments; flags add precision.
	\item \textbf{Fail gracefully.}
	When files or tools are unavailable, explain what's missing
	and proceed with remaining functionality.
	\item \textbf{Phased execution.}
	Break complex workflows into phases with user checkpoints.
	Discovery $\to$ Analysis $\to$ Proposal $\to$ Implementation $\to$ Verification.
	\item \textbf{Every skill takes iteration.}
	Expect about three weeks of refinement per skill.
	Start minimal and improve based on actual usage.
	\item \textbf{Start permissive, tighten over time.}
	For confirmation guidelines (when should Claude ask before acting?),
	start by letting Claude proceed and add guardrails as you learn
	where mistakes happen.
\end{enumerate}
